% Môn: Vật lý
% Lớp: 11
% Chương: Chương I. Dao động
% Bài: L11C1 Bài 1. Dao động điều hòa · Xác định độ lệch pha giữa hai dao động điều hòa cùng chu kì
% ===== Câu 41 =====
% ID: L11C1B1-97-TN
% Mức: NB
\dangbt{Xác định độ lệch pha giữa hai dao động điều hòa cùng chu kì}
\begin{ex}
% ID: L11C1B1-97-TN
% Mức: NB
Hai vật dao động điều hòa cùng chu kì với phương trình: $x_1 = 4\cos(2\pi t)$ cm và $x_2 = 4\cos(2\pi t + \frac{\pi}{3})$ cm. Độ lệch pha giữa hai dao động là bao nhiêu?
\choice
{$\Delta\varphi = 0$}
{$\Delta\varphi = \frac{\pi}{6}$}
{\True $\Delta\varphi = \frac{\pi}{3}$}
{$\Delta\varphi = \frac{\pi}{2}$}
\loigiai{
Độ lệch pha giữa hai dao động cùng chu kì được tính bằng hiệu pha ban đầu:
  \begin{aligned}
 $\Delta\varphi&=\varphi_2$ - $\varphi_1$ &= $\frac{\pi}{3}$ - 0 
&= $\frac{\pi}{3}\\text{rad}$.
 \end{aligned} 
 Vậy độ lệch pha là $\frac{\pi}{3}$ rad.
}
\end{ex}

% ===== Câu 42 =====
% ID: L11C1B1-98-TN
% Mức: NB
\begin{ex}
% ID: L11C1B1-98-TN
% Mức: NB
Hai vật dao động điều hòa cùng chu kì với phương trình: $x_1 = 4\cos(2\pi t)$ cm và $x_2 = 4\cos(2\pi t + \frac{\pi}{3})$ cm. Độ lệch pha giữa hai dao động là bao nhiêu?
\choice
{$\Delta\varphi = 0$}
{$\Delta\varphi = \frac{\pi}{6}$}
{\True $\Delta\varphi = \frac{\pi}{3}$}
{$\Delta\varphi = \frac{\pi}{2}$}
\loigiai{
Độ lệch pha giữa hai dao động cùng chu kì được tính bằng hiệu pha ban đầu:
  \begin{aligned}
 $\Delta\varphi&=\varphi_2$ - $\varphi_1$ &= $\frac{\pi}{3}$ - 0 
&= $\frac{\pi}{3}\\text{rad}$.
 \end{aligned} 
 Vậy độ lệch pha là $\frac{\pi}{3}$ rad.
}
\end{ex}

% ===== Câu 43 =====
% ID: L11C1B1-99-TN
% Mức: TH
\begin{ex}
% ID: L11C1B1-99-TN
% Mức: TH
Hai dao động điều hòa có phương trình: $x_1 = 6\sin(5\pi t)$ cm và $x_2 = 6\sin(5\pi t + \frac{\pi}{2})$ cm. Hai dao động có mối liên hệ gì?
\choice
{Hai dao động cùng pha}
{\True Hai dao động vuông pha}
{Hai dao động ngược pha}
{Hai dao động lệch pha $\frac{\pi}{3}$}
\loigiai{
Độ lệch pha giữa hai dao động:
  \begin{aligned}
 $\Delta\varphi&=\varphi_2$ - $\varphi_1$ &= $\frac{\pi}{2}$ - 0 
&= $\frac{\pi}{2}\\text{rad}$.
 \end{aligned} 
 Khi $\Delta\varphi = \frac{\pi}{2}$, hai dao động gọi là vuông pha.
}
\end{ex}

% ===== Câu 44 =====
% ID: L11C1B1-100-TN
% Mức: TH
\begin{ex}
% ID: L11C1B1-100-TN
% Mức: TH
Hai dao động điều hòa có phương trình: $x_1 = 6\sin(5\pi t)$ cm và $x_2 = 6\sin(5\pi t + \frac{\pi}{2})$ cm. Hai dao động có mối liên hệ gì?
\choice
{Hai dao động cùng pha}
{\True Hai dao động vuông pha}
{Hai dao động ngược pha}
{Hai dao động lệch pha $\frac{\pi}{3}$}
\loigiai{
Độ lệch pha giữa hai dao động:
  \begin{aligned}
 $\Delta\varphi&=\varphi_2$ - $\varphi_1$ &= $\frac{\pi}{2}$ - 0 
&= $\frac{\pi}{2}\\text{rad}$.
 \end{aligned} 
 Khi $\Delta\varphi = \frac{\pi}{2}$, hai dao động gọi là vuông pha.
}
\end{ex}

% ===== Câu 45 =====
% ID: L11C1B1-101-TN
% Mức: TH
\begin{ex}
% ID: L11C1B1-101-TN
% Mức: TH
Hai dao động điều hòa có phương trình: $x_1 = 6\sin(5\pi t)$ cm và $x_2 = 6\sin(5\pi t + \frac{\pi}{2})$ cm. Hai dao động có mối liên hệ gì?
\choice
{Hai dao động cùng pha}
{\True Hai dao động vuông pha}
{Hai dao động ngược pha}
{Hai dao động lệch pha $\frac{\pi}{3}$}
\loigiai{
Độ lệch pha giữa hai dao động:
  \begin{aligned}
 $\Delta\varphi&=\varphi_2$ - $\varphi_1$ &= $\frac{\pi}{2}$ - 0 
&= $\frac{\pi}{2}\\text{rad}$.
 \end{aligned} 
 Khi $\Delta\varphi = \frac{\pi}{2}$, hai dao động gọi là vuông pha.
}
\end{ex}

% ===== Câu 46 =====
% ID: L11C1B1-102-TN
% Mức: TH
\begin{ex}
% ID: L11C1B1-102-TN
% Mức: TH
Hai dao động cùng chu kì với phương trình: $x_1 = 2\cos(\pi t + \frac{\pi}{3})$ cm và $x_2 = 2\cos(\pi t + \pi)$ cm. Xác định mối liên hệ pha giữa chúng.
\choice
{\True Hai dao động lệch pha $\frac{2\pi}{3}$}
{Hai dao động cùng pha}
{Hai dao động vuông pha}
{Hai dao động ngược pha}
\loigiai{
Độ lệch pha:
  \begin{aligned}
 $\Delta\varphi&=\varphi_2$ - $\varphi_1$ &= $\pi$ - $\frac{\pi}{3}$ &= $\frac{2\pi}{3}\\text{rad}$.
 \end{aligned} 
 Vậy hai dao động lệch pha $\frac{2\pi}{3}$.
}
\end{ex}

% ===== Câu 47 =====
% ID: L11C1B1-103-TN
% Mức: TH
\begin{ex}
% ID: L11C1B1-103-TN
% Mức: TH
Hai dao động cùng chu kì với phương trình: $x_1 = 2\cos(\pi t + \frac{\pi}{3})$ cm và $x_2 = 2\cos(\pi t + \pi)$ cm. Xác định mối liên hệ pha giữa chúng.
\choice
{\True Hai dao động lệch pha $\frac{2\pi}{3}$}
{Hai dao động cùng pha}
{Hai dao động vuông pha}
{Hai dao động ngược pha}
\loigiai{
Độ lệch pha:
  \begin{aligned}
 $\Delta\varphi&=\varphi_2$ - $\varphi_1$ &= $\pi$ - $\frac{\pi}{3}$ &= $\frac{2\pi}{3}\\text{rad}$.
 \end{aligned} 
 Vậy hai dao động lệch pha $\frac{2\pi}{3}$.
}
\end{ex}

% ===== Câu 48 =====
% ID: L11C1B1-104-TN
% Mức: TH
\begin{ex}
% ID: L11C1B1-104-TN
% Mức: TH
Hai dao động cùng chu kì với phương trình: $x_1 = 2\cos(\pi t + \frac{\pi}{3})$ cm và $x_2 = 2\cos(\pi t + \pi)$ cm. Xác định mối liên hệ pha giữa chúng.
\choice
{\True Hai dao động lệch pha $\frac{2\pi}{3}$}
{Hai dao động cùng pha}
{Hai dao động vuông pha}
{Hai dao động ngược pha}
\loigiai{
Độ lệch pha:
  \begin{aligned}
 $\Delta\varphi&=\varphi_2$ - $\varphi_1$ &= $\pi$ - $\frac{\pi}{3}$ &= $\frac{2\pi}{3}\\text{rad}$.
 \end{aligned} 
 Vậy hai dao động lệch pha $\frac{2\pi}{3}$.
}
\end{ex}

% ===== Câu 49 =====
% ID: L11C1B1-105-TN
% Mức: VD
\begin{ex}
% ID: L11C1B1-105-TN
% Mức: VD
Phương trình dao động của hai vật là: $x_1 = 3\cos(4\pi t + \frac{\pi}{4})$ cm và $x_2 = 3\cos(4\pi t - \frac{3\pi}{4})$ cm. Tính độ lệch pha giữa hai dao động.
\choice
{$\Delta\varphi = 0$}
{$\Delta\varphi = \frac{\pi}{2}$}
{\True $\Delta\varphi = \pi$}
{$\Delta\varphi = 2\pi$}
\loigiai{
Độ lệch pha:
  \begin{aligned}
 $\Delta\varphi&=\varphi_2$ - $\varphi_1$ &= $(-\frac{3\pi}{4})$ - $\frac{\pi}{4}$ &= - $\pi$.
 \end{aligned} 
 Giá trị tuyệt đối của độ lệch pha là $|\Delta\varphi| = \pi$ rad. Hai dao động ngược pha.
}
\end{ex}

% ===== Câu 50 =====
% ID: L11C1B1-106-TN
% Mức: VD
\begin{ex}
% ID: L11C1B1-106-TN
% Mức: VD
Phương trình dao động của hai vật là: $x_1 = 3\cos(4\pi t + \frac{\pi}{4})$ cm và $x_2 = 3\cos(4\pi t - \frac{3\pi}{4})$ cm. Tính độ lệch pha giữa hai dao động.
\choice
{$\Delta\varphi = 0$}
{$\Delta\varphi = \frac{\pi}{2}$}
{\True $\Delta\varphi = \pi$}
{$\Delta\varphi = 2\pi$}
\loigiai{
Độ lệch pha:
  \begin{aligned}
 $\Delta\varphi&=\varphi_2$ - $\varphi_1$ &= $(-\frac{3\pi}{4})$ - $\frac{\pi}{4}$ &= - $\pi$.
 \end{aligned} 
 Giá trị tuyệt đối của độ lệch pha là $|\Delta\varphi| = \pi$ rad. Hai dao động ngược pha.
}
\end{ex}

% ===== Câu 55 =====
% ID: L11C1B1-107-TN
% Mức: VD
\begin{ex}
% ID: L11C1B1-107-TN
% Mức: VD
Phương trình dao động của hai vật là: $x_1 = 3\cos(4\pi t + \frac{\pi}{4})$ cm và $x_2 = 3\cos(4\pi t - \frac{3\pi}{4})$ cm. Tính độ lệch pha giữa hai dao động.
\choice
{$\Delta\varphi = 0$}
{$\Delta\varphi = \frac{\pi}{2}$}
{\True $\Delta\varphi = \pi$}
{$\Delta\varphi = 2\pi$}
\loigiai{
Độ lệch pha:
  \begin{aligned}
 $\Delta\varphi&=\varphi_2$ - $\varphi_1$ &= $(-\frac{3\pi}{4})$ - $\frac{\pi}{4}$ &= - $\pi$.
 \end{aligned} 
 Giá trị tuyệt đối của độ lệch pha là $|\Delta\varphi| = \pi$ rad. Hai dao động ngược pha.
}
\end{ex}

% ===== Câu 86 =====
% ID: L11C1B1-108-TLN
% Mức: TH
\begin{ex}
% ID: L11C1B1-108-TLN
% Mức: TH
Hai dao động cùng tần số có pha ban đầu $\pi/6$ và $\pi/2.$ Độ lệch pha bằng bao nhiêu lần $\pi$?
\shortans{01/03/2026}
\end{ex}

% ===== Câu 90 =====
% ID: L11C1B1-109-TLN
% Mức: VD
\begin{ex}
% ID: L11C1B1-109-TLN
% Mức: VD
Hai dao động cùng tần số có pha ban đầu $\pi/6$ và $\pi/2.$ Độ lệch pha bằng bao nhiêu lần $\pi$?
\shortans{01/03/2026}
\end{ex}

% ===== Câu 94 =====
% ID: L11C1B1-110-TLN
% Mức: VD
\begin{ex}
% ID: L11C1B1-110-TLN
% Mức: VD
Hai dao động cùng tần số có pha ban đầu $\pi/6$ và $\pi/2.$ Độ lệch pha bằng bao nhiêu lần $\pi$?
\shortans{01/03/2026}
\end{ex}

% ===== Câu 97 =====
% ID: L11C1B1-111-DS
% Mức: TH
\begin{ex}
% ID: L11C1B1-111-DS
% Mức: TH
Vật dao động theo $x=4\cos(2\pi t/1+\pi/3)$ cm. Xác định pha, li độ và chiều chuyển động tại $t=0$.
\choiceTF

\end{ex}

% ===== Câu 100 =====
% ID: L11C1B1-112-DS
% Mức: VD
\begin{ex}
% ID: L11C1B1-112-DS
% Mức: VD
Vật dao động theo $x=6\cos(2\pi t/2+\pi/3)$ cm. Xác định pha, li độ và chiều chuyển động tại $t=0$.
\choiceTF

\end{ex}

% ===== Câu 103 =====
% ID: L11C1B1-113-DS
% Mức: VD
\begin{ex}
% ID: L11C1B1-113-DS
% Mức: VD
Vật dao động theo $x=8\cos(2\pi t/4+\pi/3)$ cm. Xác định pha, li độ và chiều chuyển động tại $t=0$.
\choiceTF

\end{ex}

% ===== Câu 177 =====
% ID: L11C1B1-114-DS
% Mức: TH
\begin{ex}
% ID: L11C1B1-114-DS
% Mức: TH
Cho hai vật dao động điều hòa cùng chu kì có các phương trình li độ lần lượt là $x_1 = A_1\cos(\omega t + \varphi_1)$ và $x_2 = A_2\cos(\omega t + \varphi_2)$. Xét về mặt lí thuyết độ lệch pha giữa chúng:
\choiceTF
{\True Độ lệch pha của dao động thứ nhất so với dao động thứ hai được xác định bằng biểu thức $\Delta\varphi = \varphi_1 - \varphi_2$.}
{Nếu $\Delta\varphi = 2k\pi$ với $k$ là số nguyên, ta kết luận hai dao động này ngược pha nhau.}
{\True Độ lệch pha giữa hai dao động cùng chu kì luôn là một hằng số không thay đổi theo thời gian.}
{\True Nếu dao động thứ nhất trễ pha hơn dao động thứ hai một góc bằng $\pi$ rad thì hai dao động ngược pha nhau.}
\loigiai{
\begin{itemchoice}
\item (Đúng) Độ lệch pha của dao động thứ nhất so với dao động thứ hai được xác định bằng biểu thức $\Delta\varphi = \varphi_1 - \varphi_2$. (Vì): ộ lệch pha của dao động thứ nhất so với dao động thứ hai được xác định bằng hiệu số pha của chúng, tức là $\Delta\varphi = \varphi_1 - \varphi_2$
\item (Sai) Nếu $\Delta\varphi = 2k\pi$ với $k$ là số nguyên, ta kết luận hai dao động này ngược pha nhau. (Vì): Nếu độ lệch pha $\Delta\varphi = 2k\pi$ với $k$ là số nguyên, hai dao động này cùng pha với nhau, không phải ngược pha
\item (Đúng) Độ lệch pha giữa hai dao động cùng chu kì luôn là một hằng số không thay đổi theo thời gian. (Vì): ộ lệch pha giữa hai dao động cùng chu kì là $\Delta\varphi = (\omega t + \varphi_1) - (\omega t + \varphi_2) = \varphi_1 - \varphi_2$, là một hằng số không phụ thuộc vào thời gian $t$
\item (Đúng) Nếu dao động thứ nhất trễ pha hơn dao động thứ hai một góc bằng $\pi$ rad thì hai dao động ngược pha nhau. (Vì): Nếu dao động thứ nhất trễ pha hơn dao động thứ hai một góc bằng $\pi$ rad, tức là $\varphi_1 - \varphi_2 = -\pi$ hoặc $\varphi_2 - \varphi_1 = \pi$, thì hai dao động này ngược pha nhau
\end{itemchoice}
}
\end{ex}

% ===== Câu 179 =====
% ID: L11C1B1-115-DS
% Mức: TH
\begin{ex}
% ID: L11C1B1-115-DS
% Mức: TH
Cho hai vật dao động điều hòa cùng chu kì có các phương trình li độ lần lượt là $x_1 = A_1\cos(\omega t + \varphi_1)$ và $x_2 = A_2\cos(\omega t + \varphi_2)$. Xét về mặt lí thuyết độ lệch pha giữa chúng:
\choiceTF
{\True Độ lệch pha của dao động thứ nhất so với dao động thứ hai được xác định bằng biểu thức $\Delta\varphi = \varphi_1 - \varphi_2$.}
{Nếu $\Delta\varphi = 2k\pi$ với $k$ là số nguyên, ta kết luận hai dao động này ngược pha nhau.}
{\True Độ lệch pha giữa hai dao động cùng chu kì luôn là một hằng số không thay đổi theo thời gian.}
{\True Nếu dao động thứ nhất trễ pha hơn dao động thứ hai một góc bằng $\pi$ rad thì hai dao động ngược pha nhau.}
\loigiai{
\begin{itemchoice}
\item (Đúng) Độ lệch pha của dao động thứ nhất so với dao động thứ hai được xác định bằng biểu thức $\Delta\varphi = \varphi_1 - \varphi_2$. (Vì): ộ lệch pha của dao động thứ nhất so với dao động thứ hai được xác định bằng hiệu số pha của chúng, tức là $\Delta\varphi = \varphi_1 - \varphi_2$
\item (Sai) Nếu $\Delta\varphi = 2k\pi$ với $k$ là số nguyên, ta kết luận hai dao động này ngược pha nhau. (Vì): Nếu độ lệch pha $\Delta\varphi = 2k\pi$ với $k$ là số nguyên, hai dao động này cùng pha với nhau, không phải ngược pha
\item (Đúng) Độ lệch pha giữa hai dao động cùng chu kì luôn là một hằng số không thay đổi theo thời gian. (Vì): ộ lệch pha giữa hai dao động cùng chu kì là $\Delta\varphi = (\omega t + \varphi_1) - (\omega t + \varphi_2) = \varphi_1 - \varphi_2$, là một hằng số không phụ thuộc vào thời gian $t$
\item (Đúng) Nếu dao động thứ nhất trễ pha hơn dao động thứ hai một góc bằng $\pi$ rad thì hai dao động ngược pha nhau. (Vì): Nếu dao động thứ nhất trễ pha hơn dao động thứ hai một góc bằng $\pi$ rad, tức là $\varphi_1 - \varphi_2 = -\pi$ hoặc $\varphi_2 - \varphi_1 = \pi$, thì hai dao động này ngược pha nhau
\end{itemchoice}
}
\end{ex}

% ===== Câu 181 =====
% ID: L11C1B1-116-DS
% Mức: TH
\begin{ex}
% ID: L11C1B1-116-DS
% Mức: TH
Cho hai vật dao động điều hòa cùng chu kì có các phương trình li độ lần lượt là $x_1 = A_1\cos(\omega t + \varphi_1)$ và $x_2 = A_2\cos(\omega t + \varphi_2)$. Xét về mặt lí thuyết độ lệch pha giữa chúng:
\choiceTF
{\True Độ lệch pha của dao động thứ nhất so với dao động thứ hai được xác định bằng biểu thức $\Delta\varphi = \varphi_1 - \varphi_2$.}
{Nếu $\Delta\varphi = 2k\pi$ với $k$ là số nguyên, ta kết luận hai dao động này ngược pha nhau.}
{\True Độ lệch pha giữa hai dao động cùng chu kì luôn là một hằng số không thay đổi theo thời gian.}
{\True Nếu dao động thứ nhất trễ pha hơn dao động thứ hai một góc bằng $\pi$ rad thì hai dao động ngược pha nhau.}
\loigiai{
\begin{itemchoice}
\item (Đúng) Độ lệch pha của dao động thứ nhất so với dao động thứ hai được xác định bằng biểu thức $\Delta\varphi = \varphi_1 - \varphi_2$. (Vì): ộ lệch pha của dao động thứ nhất so với dao động thứ hai được xác định bằng hiệu số pha của chúng, tức là $\Delta\varphi = \varphi_1 - \varphi_2$
\item (Sai) Nếu $\Delta\varphi = 2k\pi$ với $k$ là số nguyên, ta kết luận hai dao động này ngược pha nhau. (Vì): Nếu độ lệch pha $\Delta\varphi = 2k\pi$ với $k$ là số nguyên, hai dao động này cùng pha với nhau, không phải ngược pha
\item (Đúng) Độ lệch pha giữa hai dao động cùng chu kì luôn là một hằng số không thay đổi theo thời gian. (Vì): ộ lệch pha giữa hai dao động cùng chu kì là $\Delta\varphi = (\omega t + \varphi_1) - (\omega t + \varphi_2) = \varphi_1 - \varphi_2$, là một hằng số không phụ thuộc vào thời gian $t$
\item (Đúng) Nếu dao động thứ nhất trễ pha hơn dao động thứ hai một góc bằng $\pi$ rad thì hai dao động ngược pha nhau. (Vì): Nếu dao động thứ nhất trễ pha hơn dao động thứ hai một góc bằng $\pi$ rad, tức là $\varphi_1 - \varphi_2 = -\pi$ hoặc $\varphi_2 - \varphi_1 = \pi$, thì hai dao động này ngược pha nhau
\end{itemchoice}
}
\end{ex}

% ===== Câu 183 =====
% ID: L11C1B1-117-DS
% Mức: TH
\begin{ex}
% ID: L11C1B1-117-DS
% Mức: TH
Hai con lắc đơn thực hiện dao động điều hòa tại cùng một nơi với cùng chu kì dao động. Quan sát đồ thị dao động của chúng, người ta nhận thấy tại mọi thời điểm, khi vật nặng của con lắc thứ nhất đạt li độ cực đại dương thì vật nặng của con lắc thứ hai đạt li độ cực đại âm. Khảo sát tính đúng/sai của các nhận định sau:
\choiceTF
{Dao động của hai con lắc ở trạng thái ngược pha với nhau.}
{\True Độ lệch pha giữa hai dao động của hai con lắc có giá trị bằng một số lẻ lần $\pi$ rad.}
{Chu kì dao động của hai con lắc này bắt buộc phải khác nhau để duy trì trạng thái lệch pha đó.}
{Tại mọi thời điểm trong quá trình dao động, nếu con lắc thứ nhất có li độ dương thì con lắc thứ hai có li độ âm.}
\loigiai{
\begin{itemchoice}
\item (Sai) Dao động của hai con lắc ở trạng thái ngược pha với nhau. (Vì): (Sai) Dao động của hai con lắc ở trạng thái ngược pha với nhau
\item (Đúng) Độ lệch pha giữa hai dao động của hai con lắc có giá trị bằng một số lẻ lần $\pi$ rad. (Vì): Hai dao động ngược pha có độ lệch pha thỏa mãn biểu thức $\Delta \varphi = (2k+1)\pi$ rad, với $k$ là một số nguyên
\item (Sai) Chu kì dao động của hai con lắc này bắt buộc phải khác nhau để duy trì trạng thái lệch pha đó. (Vì): Để hai dao động duy trì trạng thái lệch pha không đổi theo thời gian, chúng bắt buộc phải có cùng chu kì và cùng tần số
\item (Sai) Tại mọi thời điểm trong quá trình dao động, nếu con lắc thứ nhất có li độ dương thì con lắc thứ hai có li độ âm. (Vì): (Sai) Tại mọi thời điểm trong quá trình dao động, nếu con lắc thứ nhất có li độ dương thì con lắc thứ hai có li độ âm
\end{itemchoice}
}
\end{ex}

% ===== Câu 185 =====
% ID: L11C1B1-118-DS
% Mức: TH
\begin{ex}
% ID: L11C1B1-118-DS
% Mức: TH
Hai con lắc đơn thực hiện dao động điều hòa tại cùng một nơi với cùng chu kì dao động. Quan sát đồ thị dao động của chúng, người ta nhận thấy tại mọi thời điểm, khi vật nặng của con lắc thứ nhất đạt li độ cực đại dương thì vật nặng của con lắc thứ hai đạt li độ cực đại âm. Khảo sát tính đúng/sai của các nhận định sau:
\choiceTF
{Dao động của hai con lắc ở trạng thái ngược pha với nhau.}
{\True Độ lệch pha giữa hai dao động của hai con lắc có giá trị bằng một số lẻ lần $\pi$ rad.}
{Chu kì dao động của hai con lắc này bắt buộc phải khác nhau để duy trì trạng thái lệch pha đó.}
{Tại mọi thời điểm trong quá trình dao động, nếu con lắc thứ nhất có li độ dương thì con lắc thứ hai có li độ âm.}
\loigiai{
\begin{itemchoice}
\item (Sai) Dao động của hai con lắc ở trạng thái ngược pha với nhau. (Vì): (Sai) Dao động của hai con lắc ở trạng thái ngược pha với nhau
\item (Đúng) Độ lệch pha giữa hai dao động của hai con lắc có giá trị bằng một số lẻ lần $\pi$ rad. (Vì): Hai dao động ngược pha có độ lệch pha thỏa mãn biểu thức $\Delta \varphi = (2k+1)\pi$ rad, với $k$ là một số nguyên
\item (Sai) Chu kì dao động của hai con lắc này bắt buộc phải khác nhau để duy trì trạng thái lệch pha đó. (Vì): Để hai dao động duy trì trạng thái lệch pha không đổi theo thời gian, chúng bắt buộc phải có cùng chu kì và cùng tần số
\item (Sai) Tại mọi thời điểm trong quá trình dao động, nếu con lắc thứ nhất có li độ dương thì con lắc thứ hai có li độ âm. (Vì): (Sai) Tại mọi thời điểm trong quá trình dao động, nếu con lắc thứ nhất có li độ dương thì con lắc thứ hai có li độ âm
\end{itemchoice}
}
\end{ex}

% ===== Câu 187 =====
% ID: L11C1B1-119-DS
% Mức: TH
\begin{ex}
% ID: L11C1B1-119-DS
% Mức: TH
Hai con lắc đơn thực hiện dao động điều hòa tại cùng một nơi với cùng chu kì dao động. Quan sát đồ thị dao động của chúng, người ta nhận thấy tại mọi thời điểm, khi vật nặng của con lắc thứ nhất đạt li độ cực đại dương thì vật nặng của con lắc thứ hai đạt li độ cực đại âm. Khảo sát tính đúng/sai của các nhận định sau:
\choiceTF
{Dao động của hai con lắc ở trạng thái ngược pha với nhau.}
{\True Độ lệch pha giữa hai dao động của hai con lắc có giá trị bằng một số lẻ lần $\pi$ rad.}
{Chu kì dao động của hai con lắc này bắt buộc phải khác nhau để duy trì trạng thái lệch pha đó.}
{Tại mọi thời điểm trong quá trình dao động, nếu con lắc thứ nhất có li độ dương thì con lắc thứ hai có li độ âm.}
\loigiai{
\begin{itemchoice}
\item (Sai) Dao động của hai con lắc ở trạng thái ngược pha với nhau. (Vì): (Sai) Dao động của hai con lắc ở trạng thái ngược pha với nhau
\item (Đúng) Độ lệch pha giữa hai dao động của hai con lắc có giá trị bằng một số lẻ lần $\pi$ rad. (Vì): Hai dao động ngược pha có độ lệch pha thỏa mãn biểu thức $\Delta \varphi = (2k+1)\pi$ rad, với $k$ là một số nguyên
\item (Sai) Chu kì dao động của hai con lắc này bắt buộc phải khác nhau để duy trì trạng thái lệch pha đó. (Vì): Để hai dao động duy trì trạng thái lệch pha không đổi theo thời gian, chúng bắt buộc phải có cùng chu kì và cùng tần số
\item (Sai) Tại mọi thời điểm trong quá trình dao động, nếu con lắc thứ nhất có li độ dương thì con lắc thứ hai có li độ âm. (Vì): (Sai) Tại mọi thời điểm trong quá trình dao động, nếu con lắc thứ nhất có li độ dương thì con lắc thứ hai có li độ âm
\end{itemchoice}
}
\end{ex}

% ===== Câu 189 =====
% ID: L11C1B1-120-DS
% Mức: VD
\begin{ex}
% ID: L11C1B1-120-DS
% Mức: VD
Hai vật dao động điều hòa cùng tần số dọc theo hai đường thẳng song song cạnh nhau với phương trình li độ lần lượt là
	$x_1=5\cos\left(10\pi t+\dfrac{\pi}{3}\right)\;(\mathrm{cm})$
	và
	$x_2=10\cos\left(10\pi t-\dfrac{\pi}{6}\right)\;(\mathrm{cm}).$
	Khảo sát mối quan hệ pha giữa hai dao động:
\choiceTF
{\True Dao động thứ nhất sớm pha hơn dao động thứ hai một góc bằng $\dfrac{\pi}{2}$ rad.}
{\True Hai dao động này ở trạng thái vuông pha với nhau tại mọi thời điểm.}
{\True Biên độ của dao động thứ hai lớn gấp hai lần biên độ của dao động thứ nhất.}
{Tại thời điểm ban đầu $t=0$, cả hai vật đều đang chuyển động theo chiều dương.}
\loigiai{
\begin{itemchoice}
\item (Đúng) Dao động thứ nhất sớm pha hơn dao động thứ hai một góc bằng $\dfrac{\pi}{2}$ rad. (Vì): ộ lệch pha giữa hai dao động là
\item (Đúng) Hai dao động này ở trạng thái vuông pha với nhau tại mọi thời điểm. (Vì): Hai dao động có cùng tần số góc nên độ lệch pha giữa chúng không đổi theo thời gian. Do
\item (Đúng) Biên độ của dao động thứ hai lớn gấp hai lần biên độ của dao động thứ nhất. (Vì): Biên độ của hai dao động lần lượt là
\item (Sai) Tại thời điểm ban đầu $t=0$, cả hai vật đều đang chuyển động theo chiều dương. (Vì): Vận tốc của hai vật là
\end{itemchoice}
}
\end{ex}

% ===== Câu 191 =====
% ID: L11C1B1-121-DS
% Mức: VD
\begin{ex}
% ID: L11C1B1-121-DS
% Mức: VD
Hai vật dao động điều hòa cùng tần số dọc theo hai đường thẳng song song cạnh nhau với phương trình li độ lần lượt là
	$x_1=5\cos\left(10\pi t+\dfrac{\pi}{3}\right)\;(\mathrm{cm})$
	và
	$x_2=10\cos\left(10\pi t-\dfrac{\pi}{6}\right)\;(\mathrm{cm}).$
	Khảo sát mối quan hệ pha giữa hai dao động:
\choiceTF
{\True Dao động thứ nhất sớm pha hơn dao động thứ hai một góc bằng $\dfrac{\pi}{2}$ rad.}
{\True Hai dao động này ở trạng thái vuông pha với nhau tại mọi thời điểm.}
{\True Biên độ của dao động thứ hai lớn gấp hai lần biên độ của dao động thứ nhất.}
{Tại thời điểm ban đầu $t=0$, cả hai vật đều đang chuyển động theo chiều dương.}
\loigiai{
\begin{itemchoice}
\item (Đúng) Dao động thứ nhất sớm pha hơn dao động thứ hai một góc bằng $\dfrac{\pi}{2}$ rad. (Vì): ộ lệch pha giữa hai dao động là
\item (Đúng) Hai dao động này ở trạng thái vuông pha với nhau tại mọi thời điểm. (Vì): Hai dao động có cùng tần số góc nên độ lệch pha giữa chúng không đổi theo thời gian. Do
\item (Đúng) Biên độ của dao động thứ hai lớn gấp hai lần biên độ của dao động thứ nhất. (Vì): Biên độ của hai dao động lần lượt là
\item (Sai) Tại thời điểm ban đầu $t=0$, cả hai vật đều đang chuyển động theo chiều dương. (Vì): Vận tốc của hai vật là
\end{itemchoice}
}
\end{ex}

% ===== Câu 193 =====
% ID: L11C1B1-122-DS
% Mức: VD
\begin{ex}
% ID: L11C1B1-122-DS
% Mức: VD
Hai vật dao động điều hòa cùng tần số dọc theo hai đường thẳng song song cạnh nhau với phương trình li độ lần lượt là
	$x_1=5\cos\left(10\pi t+\dfrac{\pi}{3}\right)\;(\mathrm{cm})$
	và
	$x_2=10\cos\left(10\pi t-\dfrac{\pi}{6}\right)\;(\mathrm{cm}).$
	Khảo sát mối quan hệ pha giữa hai dao động:
\choiceTF
{\True Dao động thứ nhất sớm pha hơn dao động thứ hai một góc bằng $\dfrac{\pi}{2}$ rad.}
{\True Hai dao động này ở trạng thái vuông pha với nhau tại mọi thời điểm.}
{\True Biên độ của dao động thứ hai lớn gấp hai lần biên độ của dao động thứ nhất.}
{Tại thời điểm ban đầu $t=0$, cả hai vật đều đang chuyển động theo chiều dương.}
\loigiai{
\begin{itemchoice}
\item (Đúng) Dao động thứ nhất sớm pha hơn dao động thứ hai một góc bằng $\dfrac{\pi}{2}$ rad. (Vì): ộ lệch pha giữa hai dao động là
\item (Đúng) Hai dao động này ở trạng thái vuông pha với nhau tại mọi thời điểm. (Vì): Hai dao động có cùng tần số góc nên độ lệch pha giữa chúng không đổi theo thời gian. Do
\item (Đúng) Biên độ của dao động thứ hai lớn gấp hai lần biên độ của dao động thứ nhất. (Vì): Biên độ của hai dao động lần lượt là
\item (Sai) Tại thời điểm ban đầu $t=0$, cả hai vật đều đang chuyển động theo chiều dương. (Vì): Vận tốc của hai vật là
\end{itemchoice}
}
\end{ex}

% ===== Câu 218 =====
% ID: L11C1B1-123-TN
% Mức: NB
\begin{ex}
% ID: L11C1B1-123-TN
% Mức: NB
Hai vật dao động điều hòa với $x_1 = 5\sin(3\pi t)$ cm và $x_2 = 5\sin(3\pi t - \frac{\pi}{6})$ cm. Độ lệch pha là bao nhiêu?
\choice
{$\Delta\varphi = \frac{\pi}{6}$}
{\True $\Delta\varphi = \frac{\pi}{6}$}
{$\Delta\varphi = \frac{\pi}{3}$}
{$\Delta\varphi = \frac{\pi}{2}$}
\loigiai{
Độ lệch pha:
  \begin{aligned}
 $\Delta\varphi&=\varphi_2$ - $\varphi_1$ &= $(-\frac{\pi}{6})$ - 0 
&= - $\frac{\pi}{6}$.
 \end{aligned} 
 Giá trị tuyệt đối là $|\Delta\varphi| = \frac{\pi}{6}$ rad.
}
\end{ex}

% ===== Câu 219 =====
% ID: L11C1B1-124-TN
% Mức: NB
\begin{ex}
% ID: L11C1B1-124-TN
% Mức: NB
Hai vật dao động điều hòa với $x_1 = 5\sin(3\pi t)$ cm và $x_2 = 5\sin(3\pi t - \frac{\pi}{6})$ cm. Độ lệch pha là bao nhiêu?
\choice
{$\Delta\varphi = \frac{\pi}{6}$}
{\True $\Delta\varphi = \frac{\pi}{6}$}
{$\Delta\varphi = \frac{\pi}{3}$}
{$\Delta\varphi = \frac{\pi}{2}$}
\loigiai{
Độ lệch pha:
  \begin{aligned}
 $\Delta\varphi&=\varphi_2$ - $\varphi_1$ &= $(-\frac{\pi}{6})$ - 0 
&= - $\frac{\pi}{6}$.
 \end{aligned} 
 Giá trị tuyệt đối là $|\Delta\varphi| = \frac{\pi}{6}$ rad.
}
\end{ex}

% ===== Câu 220 =====
% ID: L11C1B1-125-TN
% Mức: NB
\begin{ex}
% ID: L11C1B1-125-TN
% Mức: NB
Hai vật dao động điều hòa với $x_1 = 5\sin(3\pi t)$ cm và $x_2 = 5\sin(3\pi t - \frac{\pi}{6})$ cm. Độ lệch pha là bao nhiêu?
\choice
{$\Delta\varphi = \frac{\pi}{6}$}
{\True $\Delta\varphi = \frac{\pi}{6}$}
{$\Delta\varphi = \frac{\pi}{3}$}
{$\Delta\varphi = \frac{\pi}{2}$}
\loigiai{
Độ lệch pha:
  \begin{aligned}
 $\Delta\varphi&=\varphi_2$ - $\varphi_1$ &= $(-\frac{\pi}{6})$ - 0 
&= - $\frac{\pi}{6}$.
 \end{aligned} 
 Giá trị tuyệt đối là $|\Delta\varphi| = \frac{\pi}{6}$ rad.
}
\end{ex}

% ===== Câu 221 =====
% ID: L11C1B1-126-TN
% Mức: NB
\begin{ex}
% ID: L11C1B1-126-TN
% Mức: NB
Hai vật dao động điều hòa cùng chu kỳ với phương trình $x_1=6\cos(2\pi t)$ cm và $x_2=6\cos(2\pi t+\dfrac{\pi}{2})$ cm. Độ lệch pha giữa hai dao động là bao nhiêu?
\choice
{$\Delta\varphi=0$}
{$\Delta\varphi=\dfrac{\pi}{4}$}
{\True $\Delta\varphi=\dfrac{\pi}{2}$}
{$\Delta\varphi=\pi$}
\loigiai{
Phương trình của vật thứ nhất: $x_1=6\cos(2\pi t)$, pha tại thời điểm $t$ là $\varphi_1=2\pi t$.
 Phương trình của vật thứ hai: $x_2=6\cos(2\pi t+\dfrac{\pi}{2})$, pha tại thời điểm $t$ là $\varphi_2=2\pi t+\dfrac{\pi}{2}$.
 Độ lệch pha giữa hai dao động:
  \begin{aligned}
 $\Delta\varphi$ &= $\varphi_2-\varphi_1$ &= $(2\pi$ t+ $\dfrac{\pi}{2})-2\pi$ t 
&= $\dfrac{\pi}{2}\\text{rad}$.
 \end{aligned} 
 Chọn đáp án tương ứng.
}
\end{ex}

% ===== Câu 222 =====
% ID: L11C1B1-127-TN
% Mức: NB
\begin{ex}
% ID: L11C1B1-127-TN
% Mức: NB
Hai vật dao động điều hòa cùng chu kỳ với phương trình $x_1=6\cos(2\pi t)$ cm và $x_2=6\cos(2\pi t+\dfrac{\pi}{2})$ cm. Độ lệch pha giữa hai dao động là bao nhiêu?
\choice
{$\Delta\varphi=0$}
{$\Delta\varphi=\dfrac{\pi}{4}$}
{\True $\Delta\varphi=\dfrac{\pi}{2}$}
{$\Delta\varphi=\pi$}
\loigiai{
Phương trình của vật thứ nhất: $x_1=6\cos(2\pi t)$, pha tại thời điểm $t$ là $\varphi_1=2\pi t$.
 Phương trình của vật thứ hai: $x_2=6\cos(2\pi t+\dfrac{\pi}{2})$, pha tại thời điểm $t$ là $\varphi_2=2\pi t+\dfrac{\pi}{2}$.
 Độ lệch pha giữa hai dao động:
  \begin{aligned}
 $\Delta\varphi$ &= $\varphi_2-\varphi_1$ &= $(2\pi$ t+ $\dfrac{\pi}{2})-2\pi$ t 
&= $\dfrac{\pi}{2}\\text{rad}$.
 \end{aligned} 
 Chọn đáp án tương ứng.
}
\end{ex}

% ===== Câu 223 =====
% ID: L11C1B1-128-TN
% Mức: NB
\begin{ex}
% ID: L11C1B1-128-TN
% Mức: NB
Hai vật dao động điều hòa cùng chu kỳ với phương trình $x_1=6\cos(2\pi t)$ cm và $x_2=6\cos(2\pi t+\dfrac{\pi}{2})$ cm. Độ lệch pha giữa hai dao động là bao nhiêu?
\choice
{$\Delta\varphi=0$}
{$\Delta\varphi=\dfrac{\pi}{4}$}
{\True $\Delta\varphi=\dfrac{\pi}{2}$}
{$\Delta\varphi=\pi$}
\loigiai{
Phương trình của vật thứ nhất: $x_1=6\cos(2\pi t)$, pha tại thời điểm $t$ là $\varphi_1=2\pi t$.
 Phương trình của vật thứ hai: $x_2=6\cos(2\pi t+\dfrac{\pi}{2})$, pha tại thời điểm $t$ là $\varphi_2=2\pi t+\dfrac{\pi}{2}$.
 Độ lệch pha giữa hai dao động:
  \begin{aligned}
 $\Delta\varphi$ &= $\varphi_2-\varphi_1$ &= $(2\pi$ t+ $\dfrac{\pi}{2})-2\pi$ t 
&= $\dfrac{\pi}{2}\\text{rad}$.
 \end{aligned} 
 Chọn đáp án tương ứng.
}
\end{ex}

% ===== Câu 224 =====
% ID: L11C1B1-129-TN
% Mức: NB
\begin{ex}
% ID: L11C1B1-129-TN
% Mức: NB
Hai vật dao động điều hòa cùng chu kì với phương trình: $x_1 = 4\cos(2\pi t)$ cm và $x_2 = 4\cos(2\pi t + \frac{\pi}{3})$ cm. Độ lệch pha giữa hai dao động là bao nhiêu?
\choice
{$\Delta\varphi = 0$}
{$\Delta\varphi = \frac{\pi}{6}$}
{\True $\Delta\varphi = \frac{\pi}{3}$}
{$\Delta\varphi = \frac{\pi}{2}$}
\loigiai{
Độ lệch pha giữa hai dao động cùng chu kì được tính bằng hiệu pha ban đầu:
  \begin{aligned}
 $\Delta\varphi&=\varphi_2$ - $\varphi_1$ &= $\frac{\pi}{3}$ - 0 
&= $\frac{\pi}{3}\\text{rad}$.
 \end{aligned} 
 Vậy độ lệch pha là $\frac{\pi}{3}$ rad.
}
\end{ex}

% ===== Câu 101 =====
% ID: L11C1B1-130-DS
% Mức: TH
\begin{ex}
% ID: L11C1B1-130-DS
% Mức: TH
Hai con lắc đơn thực hiện dao động điều hòa tại cùng một nơi với cùng chu kì dao động. Quan sát đồ thị dao động của chúng, người ta nhận thấy tại mọi thời điểm, khi vật nặng của con lắc thứ nhất đạt li độ cực đại dương thì vật nặng của con lắc thứ hai đạt li độ cực đại âm. Khảo sát tính đúng/sai của các nhận định sau:
\choiceTF
{\True Dao động của hai con lắc ở trạng thái ngược pha với nhau.}
{\True Độ lệch pha giữa hai dao động của hai con lắc có giá trị bằng một số lẻ lần $\pi$ rad.}
{Chu kì dao động của hai con lắc này bắt buộc phải khác nhau để duy trì trạng thái lệch pha đó.}
{\True Tại mọi thời điểm trong quá trình dao động, nếu con lắc thứ nhất có li độ dương thì con lắc thứ hai có li độ âm.}
\loigiai{
\begin{itemchoice}
\item (Đúng) Dao động của hai con lắc ở trạng thái ngược pha với nhau. (Vì): _2$), trạng thái của chúng luôn đối nghịch hoàn toàn, tức là ngược pha
\item (Đúng) Độ lệch pha giữa hai dao động của hai con lắc có giá trị bằng một số lẻ lần $\pi$ rad. (Vì): Mệnh đề thứ hai đúng vì hai dao động ngược pha có độ lệch pha thỏa mãn: $\Delta\varphi = (2k+1)\pi$ rad
\item (Sai) Chu kì dao động của hai con lắc này bắt buộc phải khác nhau để duy trì trạng thái lệch pha đó. (Vì): Mệnh đề thứ ba sai vì để độ lệch pha giữ nguyên là hằng số không đổi theo thời gian thì hai dao động bắt buộc phải có cùng chu kì và tần số
\item (Đúng) Tại mọi thời điểm trong quá trình dao động, nếu con lắc thứ nhất có li độ dương thì con lắc thứ hai có li độ âm.
\end{itemchoice}
}
\end{ex}

% ===== Câu 102 =====
% ID: L11C1B1-131-DS
% Mức: TH
\begin{ex}
% ID: L11C1B1-131-DS
% Mức: TH
Cho hai vật dao động điều hòa cùng chu kì có các phương trình li độ lần lượt là $x_1 = A_1\cos(\omega t + \varphi_1)$ và $x_2 = A_2\cos(\omega t + \varphi_2)$. Xét về mặt lí thuyết độ lệch pha giữa chúng:
\choiceTF
{\True Độ lệch pha của dao động thứ nhất so với dao động thứ hai được xác định bằng biểu thức $\Delta\varphi = \varphi_1 - \varphi_2$.}
{Nếu $\Delta\varphi = 2k\pi$ với $k$ là số nguyên, ta kết luận hai dao động này ngược pha nhau.}
{\True Độ lệch pha giữa hai dao động cùng chu kì luôn là một hằng số không thay đổi theo thời gian.}
{\True Nếu dao động thứ nhất trễ pha hơn dao động thứ hai một góc bằng $\pi$ rad thì hai dao động ngược pha nhau.}
\loigiai{
\begin{itemchoice}
\item (Đúng) Độ lệch pha của dao động thứ nhất so với dao động thứ hai được xác định bằng biểu thức $\Delta\varphi = \varphi_1 - \varphi_2$. (Vì): Mệnh đề thứ nhất đúng vì độ lệch pha giữa hai dao động cùng tần số góc bằng hiệu số hai pha ban đầu
\item (Sai) Nếu $\Delta\varphi = 2k\pi$ với $k$ là số nguyên, ta kết luận hai dao động này ngược pha nhau. (Vì): Mệnh đề thứ hai sai vì khi $\Delta\varphi = 2k\pi$, hai dao động luôn đồng bộ về pha nên chúng cùng pha với nhau
\item (Đúng) Độ lệch pha giữa hai dao động cùng chu kì luôn là một hằng số không thay đổi theo thời gian. (Vì): Mệnh đề thứ ba đúng vì trong biểu thức pha dao động $\Phi_1 - \Phi_2 = (\omega t + \varphi_1) - (\omega t + \varphi_2) = \varphi_1 - \varphi_2$, đại lượng thời gian $t$ bị triệt tiêu nên độ lệch pha luôn không đổi
\item (Đúng) Nếu dao động thứ nhất trễ pha hơn dao động thứ hai một góc bằng $\pi$ rad thì hai dao động ngược pha nhau. (Vì): Mệnh đề thứ tư đúng vì khi lệch pha nhau một góc bằng $\pi$ rad (hoặc một số lẻ lần $\pi$), hai dao động có trạng thái đối nghịch và ngược pha nhau
\end{itemchoice}
}
\end{ex}

% ===== Câu 103 =====
% ID: L11C1B1-132-DS
% Mức: VD
\begin{ex}
% ID: L11C1B1-132-DS
% Mức: VD
Hai dao động điều hòa cùng chu kì có phương trình li độ lần lượt là $x_1 = 4\cos(2\pi t)$ (cm) và $x_2 = 8\sin(2\pi t)$ (cm). Đánh giá mối quan hệ pha giữa hai dao động này:
\choiceTF
{\True Pha ban đầu của dao động thứ hai sau khi biến đổi về dạng cosin chuẩn dương là $-\dfrac{\pi}{2}$ rad.}
{Dao động thứ nhất trễ pha hơn dao động thứ hai một góc bằng $\dfrac{\pi}{2}$ rad.}
{\True Độ lệch pha của dao động thứ nhất so với dao động thứ hai là $\dfrac{\pi}{2}$ rad.}
{Tại mọi thời điểm, tỉ số li độ giữa hai vật $\dfrac{x_1}{x_2}$ luôn giữ một giá trị không đổi.}
\loigiai{
\begin{itemchoice}
\item (Đúng) Pha ban đầu của dao động thứ hai sau khi biến đổi về dạng cosin chuẩn dương là $-\dfrac{\pi}{2}$ rad. (Vì): Mệnh đề thứ nhất đúng vì ta biến đổi: $x_2 = 8\sin(2\pi t) = 8\cos\left(2\pi t - \dfrac{\pi}{2}\right)$
\item (Sai) Dao động thứ nhất trễ pha hơn dao động thứ hai một góc bằng $\dfrac{\pi}{2}$ rad. (Vì): Mệnh đề thứ hai sai vì $\Delta\varphi = \varphi_1 - \varphi_2 = 0 - \left(-\dfrac{\pi}{2}\right) = \dfrac{\pi}{2}$ rad $\gt  0$, nghĩa là dao động 1 sớm pha hơn dao động 2
\item (Đúng) Độ lệch pha của dao động thứ nhất so với dao động thứ hai là $\dfrac{\pi}{2}$ rad. (Vì): Mệnh đề thứ ba đúng vì hiệu số pha ban đầu của chúng có độ lớn bằng $\dfrac{\pi}{2}$ rad
\item (Sai) Tại mọi thời điểm, tỉ số li độ giữa hai vật $\dfrac{x_1}{x_2}$ luôn giữ một giá trị không đổi.
\end{itemchoice}
}
\end{ex}

% ===== Câu 105 =====
% ID: L11C1B1-133-DS
% Mức: VD
\begin{ex}
% ID: L11C1B1-133-DS
% Mức: VD
Hai vật dao động điều hòa cùng tần số dọc theo hai đường thẳng song song cạnh nhau với phương trình li độ lần lượt là $x_1 = 5\cos\left(10\pi t + \dfrac{\pi}{3}\right)$ (cm) và $x_2 = 10\cos\left(10\pi t - \dfrac{\pi}{6}\right)$ (cm). Khảo sát mối quan hệ pha giữa hai dao động:
\choiceTF
{\True Dao động thứ nhất sớm pha hơn dao động thứ hai một góc bằng $\dfrac{\pi}{2}$ rad.}
{\True Hai dao động này ở trạng thái vuông pha với nhau tại mọi thời điểm.}
{\True Biên độ của dao động thứ hai lớn gấp hai lần biên độ của dao động thứ nhất.}
{Tại thời điểm ban đầu $t = 0$, cả hai vật đều đang chuyển động theo chiều dương.}
\loigiai{
\begin{itemchoice}
\item (Đúng) Dao động thứ nhất sớm pha hơn dao động thứ hai một góc bằng $\dfrac{\pi}{2}$ rad. (Vì): Mệnh đề thứ nhất đúng vì độ lệch pha của hai dao động là: $\Delta\varphi = \varphi_1 - \varphi_2 = \dfrac{\pi}{3} - \left(-\dfrac{\pi}{6}\right) = \dfrac{\pi}{2}$ rad $\gt  0$
\item (Đúng) Hai dao động này ở trạng thái vuông pha với nhau tại mọi thời điểm. (Vì): Mệnh đề thứ hai đúng vì độ lệch pha giữa hai dao động cùng tần số có giá trị không đổi và bằng $\dfrac{\pi}{2}$ rad, đây là điều kiện vuông pha
\item (Đúng) Biên độ của dao động thứ hai lớn gấp hai lần biên độ của dao động thứ nhất. (Vì): Mệnh đề thứ ba đúng vì biên độ tương ứng của hai vật là $A_2 = 10$ cm và $A_1 = 5$ cm
\item (Sai) Tại thời điểm ban đầu $t = 0$, cả hai vật đều đang chuyển động theo chiều dương. (Vì): Mệnh đề thứ tư sai vì tại $t = 0$, vật 1 có $\varphi_1 = \dfrac{\pi}{3} \gt  0$ nên chuyển động theo chiều âm ($v_1 \lt  0$); vật 2 có $\varphi_2 = -\dfrac{\pi}{6} \lt  0$ nên chuyển động theo chiều dương ($v_2 \gt  0$
\end{itemchoice}
}
\end{ex}

% ===== Câu 106 =====
% ID: L11C1B1-134-DS
% Mức: VD
\begin{ex}
% ID: L11C1B1-134-DS
% Mức: VD
Hai vật dao động điều hòa cùng chu kì quanh vị trí cân bằng có phương trình li độ lần lượt là $x_1 = 3\cos\left(5\pi t + \dfrac{\pi}{6}\right)$ (cm) và $x_2 = -3\cos\left(5\pi t + \dfrac{\pi}{6}\right)$ (cm). Khảo sát các thông số pha của hai dao động này:
\choiceTF
{Biên độ dao động của vật thứ hai có giá trị bằng $-3$ cm.}
{\True Pha ban đầu của dao động thứ hai sau khi chuẩn hóa biên độ dương bằng $-\dfrac{5\pi}{6}$ rad.}
{\True Hai dao động này ở trạng thái ngược pha nhau.}
{\True Độ lệch pha giữa hai dao động bằng $\pi$ rad.}
\loigiai{
\begin{itemchoice}
\item (Sai) Biên độ dao động của vật thứ hai có giá trị bằng $-3$ cm. (Vì): Mệnh đề thứ nhất sai vì biên độ dao động luôn dương, do đó $A_2 = 3$ cm
\item (Đúng) Pha ban đầu của dao động thứ hai sau khi chuẩn hóa biên độ dương bằng $-\dfrac{5\pi}{6}$ rad. (Vì): Mệnh đề thứ hai đúng vì chuẩn hóa: $x_2 = -3\cos\left(5\pi t + \dfrac{\pi}{6}\right) = 3\cos\left(5\pi t + \dfrac{\pi}{6} - \pi\right) = 3\cos\left(5\pi t - \dfrac{5\pi}{6}\right)$
\item (Đúng) Hai dao động này ở trạng thái ngược pha nhau. (Vì): Mệnh đề thứ ba đúng vì phương trình li độ của hai vật liên hệ với nhau qua biểu thức $x_2 = -x_1$, do đó chúng ngược pha nhau
\item (Đúng) Độ lệch pha giữa hai dao động bằng $\pi$ rad.
\end{itemchoice}
}
\end{ex}

% ===== Câu 128 =====
% ID: L11C1B1-135-TN
% Mức: NB
\begin{ex}
% ID: L11C1B1-135-TN
% Mức: NB
Hai chất điểm dao động điều hòa cùng chu kì có các phương trình li độ tương ứng lần lượt là $x_1 = A_1\cos(\omega t + \varphi_1)$ và $x_2 = A_2\cos(\omega t + \varphi_2)$. Nếu $\varphi_1 = \varphi_2$ thì phát biểu nào sau đây đúng?
\choice
{\True Hai dao động này luôn cùng pha với nhau}
{Hai dao động này luôn ngược pha với nhau}
{Hai dao động này luôn vuông pha với nhau}
{Dao động thứ nhất sớm pha hơn dao động thứ hai}
\loigiai{
• Tính toán hiệu số pha ban đầu giữa hai phương trình dao động của hai vật: $\Delta\varphi = \varphi_1 - \varphi_2$.
 
• Theo giả thiết đề bài cho $\varphi_1 = \varphi_2$ nên ta có: $\Delta\varphi = 0$.
 
• Do độ lệch pha ban đầu bằng $0$ nên tại mọi thời điểm li độ của hai vật luôn biến thiên đồng bộ với nhau về pha dao động. Ta kết luận hai dao động cùng pha.
}
\end{ex}

% ===== Câu 129 =====
% ID: L11C1B1-136-TN
% Mức: NB
\begin{ex}
% ID: L11C1B1-136-TN
% Mức: NB
Khi hiệu số pha ban đầu của hai dao động điều hòa cùng chu kì thỏa mãn điều kiện $\Delta\varphi = \varphi_1 - \varphi_2 = (2k+1)\pi$ với $k$ là số nguyên, ta kết luận hai dao động này:
\choice
{Cùng pha với nhau}
{\True Ngược pha với nhau}
{Vuông pha với nhau}
{Lệch pha nhau một góc $60^\circ$}
\loigiai{
• Xét biểu thức độ lệch pha giữa hai dao động có giá trị bằng một số lẻ lần $\pi$: $\Delta\varphi = \pm\pi, \pm3\pi, \dots$
 
• Khi hai dao động lệch pha nhau một góc bằng một số lẻ của $\pi$, tại mọi thời điểm li độ của chúng luôn trái dấu nhau (nếu cùng biên độ thì $x_1 = -x_2$).
 
• Trạng thái chuyển động của hai vật luôn đối nghịch nhau, do đó hai dao động này được gọi là ngược pha với nhau.
}
\end{ex}

% ===== Câu 130 =====
% ID: L11C1B1-137-TN
% Mức: TH
\begin{ex}
% ID: L11C1B1-137-TN
% Mức: TH
Hai vật dao động điều hòa cùng tần số với phương trình li độ lần lượt là $x_1 = 5\cos(2\pi t)$ (cm) và $x_2 = 5\sin(2\pi t)$ (cm). Phát biểu nào sau đây là đúng?
\choice
{\True Dao động thứ nhất sớm pha $\dfrac{\pi}{2}$ so với dao động thứ hai}
{Dao động thứ nhất trễ pha $\dfrac{\pi}{2}$ so với dao động thứ hai}
{Hai dao động này dao động cùng pha với nhau}
{Hai dao động này dao động ngược pha với nhau}
\loigiai{
A. Đưa phương trình dao động thứ hai từ dạng hàm sin về dạng hàm cosin chuẩn có biên độ dương bằng hệ thức lượng giác: $\sin\alpha = \cos\left(\alpha - \dfrac{\pi}{2}\right)$.
 
B. Ta viết lại phương trình dao động của vật thứ hai: $x_2 = 5\cos\left(2\pi t - \dfrac{\pi}{2}\right)$ (cm).
 
C. Xác định pha ban đầu của từng dao động: dao động thứ nhất có $\varphi_1 = 0$ rad, dao động thứ hai sau biến đổi có $\varphi_2 = -\dfrac{\pi}{2}$ rad.
 
D. Độ lệch pha giữa dao động 1 và dao động 2 là: $\Delta\varphi = \varphi_1 - \varphi_2 = 0 - \left(-\dfrac{\pi}{2}\right) = \dfrac{\pi}{2}$ rad. Do $\Delta\varphi \gt  0$ nên dao động 1 sớm pha hơn dao động 2 một góc $\dfrac{\pi}{2}$ rad.
}
\end{ex}

% ===== Câu 131 =====
% ID: L11C1B1-138-TN
% Mức: TH
\begin{ex}
% ID: L11C1B1-138-TN
% Mức: TH
Cho hai dao động điều hòa cùng phương, cùng tần số. Quan sát đồ thị dao động ta thấy tại mọi thời điểm, khi vật thứ nhất ở vị trí biên dương thì vật thứ hai đi qua vị trí cân bằng theo chiều âm. Mối quan hệ pha giữa hai dao động này là:
\choice
{Dao động 1 ngược pha với dao động 2}
{Dao động 1 cùng pha với dao động 2}
{\True Dao động 1 sớm pha $\dfrac{\pi}{2}$ so với dao động 2}
{Dao động 1 trễ pha $\dfrac{\pi}{2}$ so với dao động 2}
\loigiai{
A. Biểu diễn trạng thái của hai vật tại thời điểm bất kì lên vòng tròn lượng giác hoặc dùng hệ thức pha.
 
B. Khi vật 1 ở biên dương thì pha dao động của nó bằng $0$ (hoặc bội nguyên của $2\pi$). Khi đó vật 2 đang ở vị trí cân bằng và đi theo chiều âm, tương ứng pha dao động bằng $\dfrac{\pi}{2}$.
 
C. Hiệu số pha dao động tại thời điểm này của vật 2 so với vật 1 là $\Phi_2 - \Phi_1 = \dfrac{\pi}{2} \Rightarrow \Phi_1 = \Phi_2 - \dfrac{\pi}{2}$.
 
D. Do đó, dao động 1 trễ pha hơn dao động 2 một góc $\dfrac{\pi}{2}$ rad, tương đương dao động 1 sớm pha hơn dao động 2 một góc là $\dfrac{3\pi}{2}$ rad, hay dao động 1 sớm pha $\dfrac{\pi}{2}$ so với dao động 2 khi xét góc lệch nhỏ nhất.
}
\end{ex}

% ===== Câu 132 =====
% ID: L11C1B1-139-TN
% Mức: VD
\begin{ex}
% ID: L11C1B1-139-TN
% Mức: VD
Hai dao động điều hòa cùng chu kì có pha ban đầu lần lượt là $\varphi_1 = \dfrac{\pi}{3}$ rad và $\varphi_2 = -\dfrac{\pi}{6}$ rad. Nhận xét nào sau đây là chính xác khi nói về độ lệch pha giữa chúng?
\choice
{Dao động thứ nhất trễ pha hơn dao động thứ hai một góc $\dfrac{\pi}{6}$ rad}
{\True Dao động thứ nhất sớm pha hơn dao động thứ hai một góc $\dfrac{\pi}{2}$ rad}
{Dao động thứ nhất sớm pha hơn dao động thứ hai một góc $\dfrac{\pi}{6}$ rad}
{Dao động thứ nhất trễ pha hơn dao động thứ hai một góc $\dfrac{\pi}{2}$ rad}
\loigiai{
• Công thức tính độ lệch pha giữa dao động thứ nhất và dao động thứ hai là: $\Delta\varphi = \varphi_1 - \varphi_2$.
 
• Thay thế các giá trị pha ban đầu đề bài cung cấp vào hệ thức: $\Delta\varphi = \dfrac{\pi}{3} - \left(-\dfrac{\pi}{6}\right) = \dfrac{\pi}{3} + \dfrac{\pi}{6} = \dfrac{\pi}{2}$ rad.
 
• Vì hiệu số $\Delta\varphi = \dfrac{\pi}{2} \gt  0$ nên ta kết luận dao động thứ nhất sớm pha hơn dao động thứ hai một góc bằng $\dfrac{\pi}{2}$ rad.
}
\end{ex}

% ===== Câu 141 =====
% ID: L11C1B1-140-TLN
% Mức: VD
\begin{ex}
% ID: L11C1B1-140-TLN
% Mức: VD
Hai dao động điều hòa cùng chu kì có phương trình lần lượt là $x_1 = 4\cos(2\pi t + \frac{\pi}{3})$ cm và $x_2 = 4\cos(2\pi t - \frac{\pi}{6})$ cm. Độ lệch pha của $x_1$ so với $x_2$ bằng bao nhiêu radian?
\shortans{0.5}
\loigiai{
Độ lệch pha của $x_1$ so với $x_2$ được tính bằng: $\Delta\varphi = \varphi_1 - \varphi_2 = \frac{\pi}{3} - \left(-\frac{\pi}{6}\right) = \frac{\pi}{3} + \frac{\pi}{6} = \frac{2\pi}{6} + \frac{\pi}{6} = \frac{3\pi}{6} = \frac{\pi}{2}$ rad. Vì $\Delta\varphi = \frac{\pi}{2} \gt  0$ nên $x_1$ sớm pha hơn $x_2$ một góc $\frac{\pi}{2}$ rad. Đổi ra số thập phân: $\frac{\pi}{2} \approx 1.5708$ rad. Tuy nhiên đề hỏi giá trị bằng bao nhiêu radian theo dạng phân số $\pi$: $\Delta\varphi = \frac{\pi}{2}$ rad, hệ số của $\pi$ là $0.5$. Đáp án: $0.5$ (tức $0.5\pi$ rad).
}
\end{ex}

% ===== Câu 142 =====
% ID: L11C1B1-141-TLN
% Mức: VD
\begin{ex}
% ID: L11C1B1-141-TLN
% Mức: VD
Hai dao động điều hòa cùng chu kì $T = 0,4$ s có phương trình $x_1 = 6\cos(5\pi t + \frac{2\pi}{3})$ cm và $x_2 = 6\cos(5\pi t + \frac{\pi}{4})$ cm. Tính độ lệch pha của $x_1$ so với $x_2$. Kết quả biểu diễn dưới dạng $k\pi$ rad, hãy cho biết giá trị của $k$ (làm tròn đến hai chữ số thập phân).
\shortans{0.42}
\loigiai{
Độ lệch pha của $x_1$ so với $x_2$: $\Delta\varphi = \varphi_1 - \varphi_2 = \frac{2\pi}{3} - \frac{\pi}{4}$. Quy đồng mẫu số: $\frac{2\pi}{3} = \frac{8\pi}{12}$, $\frac{\pi}{4} = \frac{3\pi}{12}$. Vậy $\Delta\varphi = \frac{8\pi}{12} - \frac{3\pi}{12} = \frac{5\pi}{12}$ rad. Biểu diễn dưới dạng $k\pi$: $k = \frac{5}{12} \approx 0.4167 \approx 0.42$. Vì $\Delta\varphi \gt  0$ nên $x_1$ sớm pha hơn $x_2$. Đáp án: $k \approx 0.42$.
}
\end{ex}

% ===== Câu 143 =====
% ID: L11C1B1-142-TLN
% Mức: VD
\begin{ex}
% ID: L11C1B1-142-TLN
% Mức: VD
Một chất điểm thực hiện đồng thời hai dao động điều hòa cùng chu kì. Dao động thứ nhất có phương trình $x_1 = 5\cos(4\pi t - \frac{\pi}{6})$ cm và dao động thứ hai có phương trình $x_2 = 5\cos(4\pi t + \frac{5\pi}{12})$ cm. Tính độ lệch pha của $x_1$ so với $x_2$ (tính bằng rad, lấy giá trị trong khoảng $(-\pi; \pi]$).
\shortans{-0.5}
\loigiai{
Độ lệch pha của $x_1$ so với $x_2$ được tính theo công thức: $\Delta\varphi = \varphi_1 - \varphi_2$. Trong đó $\varphi_1 = -\frac{\pi}{6}$ và $\varphi_2 = \frac{5\pi}{12}$. Ta có: $\Delta\varphi = -\frac{\pi}{6} - \frac{5\pi}{12} = -\frac{2\pi}{12} - \frac{5\pi}{12} = -\frac{7\pi}{12}$. Tính giá trị số: $\Delta\varphi = -\frac{7\pi}{12} \approx -\frac{7 \times 3,14159}{12} \approx -\frac{21,991}{12} \approx -1,833$ rad. Tuy nhiên, cần kiểm tra lại: $-\frac{7\pi}{12} \approx -1,833$ rad nằm trong khoảng $(-\pi; \pi]$ nên hợp lệ. Vậy $\Delta\varphi = -\frac{7\pi}{12} \approx -1,83$ rad. Lưu ý: Đề yêu cầu đáp số là con số không ghi đơn vị, không ghi căn thức hoặc số $\pi$, nên đáp án là $-1,83$. Kiểm tra lại: $\varphi_1 - \varphi_2 = -\frac{\pi}{6} - \frac{5\pi}{12} = -\frac{2\pi}{12} - \frac{5\pi}{12} = -\frac{7\pi}{12} \approx -1,83$ rad. Vậy $x_1$ trễ pha hơn $x_2$ một lượng $\frac{7\pi}{12}$ rad.
}
\end{ex}
